\section{引言}

本文档只整理理论估算部分，目标是建立一条可重复运行的低频电场参数化计算链：形成机制 $\rightarrow$ 总等效偶极矩 $P_0$ $\rightarrow$ UEP 结果签名与轴频调制 $\rightarrow$ 观测方向/距离边界 $\rightarrow$ 噪声谱密度下的可探测性。

计算不用于复现任何真实潜艇或舰船的电场，也不包含真实平台的 ICCP 电流、阳极位置、涂层状态、轴系结构或几何模型。报告采用公开量级参数、解析模型和敏感性扫描，服务后续实验设计和数值仿真边界设定。

本文关注五类机制或干扰：
\begin{itemize}
    \item UEP：由总等效偶极矩 $P_0$ 形成的水下电势签名；
    \item 腐蚀/SACP/ICCP：形成 $P_0$ 的电化学和阴极保护机制；
    \item 轴频调制：$P_0$ 的周期调制，表现为 $f_s,2f_s,3f_s,\ldots$ 线谱；
    \item 尾流/流动诱导电场：导电水体在地磁场中运动形成的局部低频电场；
    \item 工频泄漏：实验或环境中的 50 Hz 及倍频干扰。
\end{itemize}

海水电导率、皮肤深度公式、舰船 UEP/轴频电场模型及尾流电磁场背景来自公开资料和公开论文，具体来源在参考文献与 \texttt{../ship\_efield\_theory/report/sources.md} 中列出。
